\documentclass[11pt]{amsart}


\usepackage{bm}
\usepackage{amsmath,amsfonts, amsthm, amssymb}
\usepackage[abbrev,backrefs]{amsrefs}
%\usepackage{showlabels}
\usepackage{microtype}
\usepackage{enumerate}
\usepackage{graphicx}
%\usepackage{braket}
\graphicspath{{../../figures/}}
\usepackage[onehalfspacing]{setspace}
\usepackage{nicefrac}
\usepackage{xcolor}
\usepackage{wasysym}

%%%% for widecheck command
% code from mathabx.sty and mathabx.dcl
\DeclareFontFamily{U}{mathx}{\hyphenchar\font45}
\DeclareFontShape{U}{mathx}{m}{n}{
      <5> <6> <7> <8> <9> <10>
      <10.95> <12> <14.4> <17.28> <20.74> <24.88>
      mathx10
      }{}
\DeclareSymbolFont{mathx}{U}{mathx}{m}{n}
\DeclareFontSubstitution{U}{mathx}{m}{n}
\DeclareMathAccent{\widecheck}{0}{mathx}{"71}

\setlength{\oddsidemargin}{0pt}
\setlength{\evensidemargin}{0pt}
\setlength{\textwidth}{6.0in}
\setlength{\topmargin}{0in}
\setlength{\textheight}{8.5in}

\setlength{\parindent}{0in}
\setlength{\parskip}{5px}

\usepackage{scalerel,stackengine}
\def\apeqA{\SavedStyle\sim}
\def\apeq{\setstackgap{L}{\dimexpr.5pt+1.5\LMpt}\ensurestackMath{%
  \ThisStyle{\mathrel{\Centerstack{{\apeqA} {\apeqA} {\apeqA}}}}}}

\def\bra#1{\mathinner{\langle{#1}|}}
\def\ket#1{\mathinner{|{#1}\rangle}}
\def\braket#1{\mathinner{\langle{#1}\rangle}}
\def\Braket#1{\langle{#1}\rangle}

\def\Real{{\mathbf R}}
\def\bbC{{\mathbb C}}
\def\Sphere{{\mathbf S}}
\def\Schwartz{{\mathcal{S}}}
\def\Lebesgue{{\mathcal{L}}}
\def\Hausdorff{{\mathcal{H}}}
\def\Prob{{\mathbf P}}
\def\Var{{\mathbf{Var}\,}}
\def\Expec{{\mathbb E\,}}
\def\xp{{\mathbb E\,}}
\def\RWExpec{{\mathbb E\,}}
\def\eps{{\varepsilon}}
\def\One{{\mathbf{1}}}

\def\PhaseSpace{{\Real^{2d}}}
\def\Evol{{\mathcal{E}}} \def\xop{{\mathbf{x}}}
\def\pop{{\mathbf{p}}}
\DeclareMathOperator{\Rec}{Rec}
\DeclareMathOperator{\Err}{Err}
\def\lsim{{\,\lesssim\,}}
\def\llog{{\,\lessapprox\,}}

% the Hilbert space
\def\Hilbert{{\mathcal{H}}}

\def\bbR{{\mathbb{R}}}
\def\bbZ{{\mathbb{Z}}}
\def\bbN{{\mathbb N}}
\def\bbH{{\mathbb H}}

\newcommand{\Wigner}[1]{{\mathcal{W}_{#1}}}
\newcommand{\Ft}[1]{{\widehat{#1}}}
\newcommand{\IFt}[1]{{\widecheck{#1}}}

\newcommand{\ot}{{\leftarrow}}

\newcommand{\weight}{{\mathrm{weight}}}
\newcommand{\poly}{{\mathrm{poly}}}
\newcommand{\Poly}{{\mathrm{Poly}}}

\newcommand{\diff}{\mathop{}\!\mathrm{d}}

\numberwithin{equation}{section}

\DeclareMathOperator{\Id}{Id}
\DeclareMathOperator{\ad}{ad}
\newcommand{\AD}{{\mathbf{ad}}}
\newcommand{\F}{{\mathcal{F}}}
\newcommand{\eg}{E}
\newcommand{\noset}{\varnothing}
\newcommand{\mbf}[1]{{\mathbf{#1}}}
\newcommand{\mcal}[1]{{\mathcal{#1}}}
\newcommand{\wtild}[1]{{\widetilde{#1}}}
\newcommand{\PathSet}{{\Pi}}
\newcommand{\ddt}{{\frac{d}{dt}}}

\DeclareMathOperator{\Tr}{Tr}

\DeclareMathOperator{\Op}{Op}
\DeclareMathOperator{\Rept}{Re}
\DeclareMathOperator{\Impt}{Im}
\DeclareMathOperator{\vol}{vol}
\DeclareMathOperator{\supp}{supp}
\DeclareMathOperator{\Variance}{Var}
\DeclareMathOperator{\argmax}{argmax}
\DeclareMathOperator{\trace}{tr}
\DeclareMathOperator{\tr}{tr}
\DeclareMathOperator{\Div}{div}

\DeclareMathOperator{\Arm}{arm}
\DeclareMathOperator{\Stick}{stick}

\newtheorem{theorem}{Theorem}
\numberwithin{theorem}{section}
\newtheorem{lemma}[theorem]{Lemma}
\newtheorem{definition}[theorem]{Definition}
\newtheorem{corollary}[theorem]{Corollary}
\newtheorem{proposition}[theorem]{Proposition}
\newtheorem{example}{Example}
\newtheorem{remark}[theorem]{Remark}

\newtheorem{hope}[theorem]{Potentially Provable Lemma}
\newtheorem{assumption}{Useful Assumption}
\hfuzz=20pt

%%%%%%%%

\newcommand{\calD}{\mathcal{D}}
\newcommand{\calP}{\mathcal{P}}
\newcommand{\calS}{\mathcal{S}}
\newcommand{\calY}{\mathcal{Y}}


\newcommand{\nf}{\nicefrac}

\newcommand{\AB}[1]{\textcolor{red}{[AB: #1]}}

\newcommand{\les}{\lesssim}



\title{2-loop RBM analysis}

\begin{document}
\maketitle
The goal of this note is to explain how [ER] estimate the averaged 2-loop observable.

Let 
\begin{align*}
\calS[A]_{ab} := \delta_{ab} \sum_{c} S_{ac} A_{cc}
\end{align*}
and for later use we note that
\begin{align*}
A \calS[B] C = \sum_{d} B_{dd} A S^{d} C
\end{align*}

Let
\begin{align*}
G_{[1,2],t} = G_t S^{x_1} G_t^*,
\end{align*}
which we think of as a matrix-valued function of $x_1$.
Its deterministic approximation is given by
\begin{align*}
(M_{[1,2],t}(x_1))_{ab} = \delta_{ab}\Theta_{ax_1},
\end{align*}
which is another matrix-valued function, this time only diagonal and given in terms of the diffusion profile $\Theta$.

The diffusion profile acts linearly on functions (matrix-valued or otherwise) separately in each coordinate via
\begin{align*}
\Theta^{(j)}[f](\bs x) = \sum_{q} \Theta_{x_j q} f(x_1,\ldots,x_{j-1},q,x_{j+1},\ldots x_k)
\end{align*}

The averaged 2-loop is 
\begin{align*}
\calX^{2}_t(x_1,x_2) = \Tr[(G_{[1,2],t} - M_{[1,2],t}) S^{x_{2}})].
\end{align*}

For given imaginary part $\eta_s$, let $\Upsilon_s$ be the control function $(\Upsilon_s)_{ab} = \frac{1}{\ell_s \eta_s}\chi(|a-b|\leq \ell_s)$, where $\ell_s$ is the diffusive length-scale $W/\sqrt{\eta_s} $.
The point is that $\Tr[M_{[1,2],t}(x_1)\calS^{x_2}]\les \Upsilon_{x_1x_2}$, so the goal is to show that
$\calX^{2}$ is $\ell \eta$ better, i.e.
\begin{align*}
|\calX^{2}(x_1,x_2)| \prec (\ell \eta)^{-1} \Upsilon_{x_1x_2}.
\end{align*}
Really what we want to understand is why if we assume such a bound at some $s$ then the equation propagates it to time $t>s$. 

First, let's record the SDE satisfied by $\calX^{2}$
\begin{proposition}
We have that $\calX^{2}$ solves
\begin{align*}
&d\calX^{2} = (\Theta^{(1)}[\calX^{2}] + \Theta^{(2)}[\calX^{2]}])\:dt+ \calF + \calM,\\
&\begin{aligned}
\calF &=   - \sum_{q} \left( G_{[1,2],t }(x_1) - M_{[1,2],t}(x) \right)_{qq} \Tr[\left( G_{[1,2],t}(q) - M_{[1,2],t}(q) \right)S^{x_2}]\\
&\quad + \left(G_t (\calS[G_t] - m) G_t   \right) S^{x_1} G_t^{*} + c.c.\\
\calM &=  - G_t \: dB_t \: G_t S^{x_1} G_t^{*} + c.c.
\end{aligned}
\end{align*}
\end{proposition}
\begin{proof}
We first compute
\begin{align*}
dG_t = G_t (\calS[G_t] - m) G_t\:dt - G_t\:dB_t\:G_t
\end{align*}

We compute further that
\begin{align*}
dG_{[1,2],t} &= G_t \calS[G_t S^{x_1} G_t^{*}] G_t^{*}\: dt\\
&\quad + \left(G_t (\calS[G_t] - m) G_t   \right) S^{x_1} G_t^{*} + c.c.\\
&\quad - G_t \: dB_t \: G_t S^{x_1} G_t^{*} + c.c.
\end{align*}
Note that
\begin{align*}
(G_t \calS[G_t S^{x_1} G_t^{*}] G_t^{*})_{ab} 
&= \sum_{q} \calS[G_t S^{x} G_t^{*}]_{qq}  (G_t)_{aq} (G^{*}_{t})_{qb}\\
&= \sum_q \sum_{q'}S_{qq'} (G_t S^{x} G_t^{*})_{qq'} (G_t)_{aq} (G^{*}_{t})_{qb}\\
&= \sum_{q'} (G_{[1,2],t}(x))_{q'q'} (G_{[1,2],t}(q'))_{ab}
\end{align*}
and therefore,
\begin{align*}
d(G_{[1,2]t}(x_1) - M_{[1,2],t}(x_1)) 
&= \sum_{q} \left( G_{[1,2],t }(x_1) - M_{[1,2],t}(x_1) \right)_{qq} M_{[1,2],t}(q)\\
&\quad +\sum_q \left( M_{[1,2],t}(x_1)\right)_{qq} \left( G_{[1,2],t}(q) - M_{[1,2],t}(q) \right) \\
&\quad - \sum_{q} \left( G_{[1,2],t }(x_1) - M_{[1,2],t}(x) \right)_{qq} \left( G_{[1,2],t}(q) - M_{[1,2],t}(q) \right)\\
&\quad + \left(G_t (\calS[G_t] - m) G_t   \right) S^{x_1} G_t^{*} + c.c.\\
&\quad - G_t \: dB_t \: G_t S^{x_1} G_t^{*} + c.c.
\end{align*}
Note that
\begin{align*}
\Tr\left[ M_{[1,2],t}(q)S^{x_2} \right]  
&= \sum_{q'} \Theta_{qq'}S_{q'x_2}\\
&= (\Theta S)_{qx_2}\\
&= (S \Theta)_{qx_2}\\
&= \sum_{q'}S_{qq'} \Theta_{q'x_2},
\end{align*}
since $\Theta$ commutes with $S$ by definition.

Multiplying by $S^{x_2}$ and taking traces, we find by the above that
\begin{align*}
\Tr(\sum_{q} \left( G_{[1,2],t }(x_1) - M_{[1,2],t}(x_1) \right)_{qq} M_{[1,2],t}(q) S^{x_2})
=  \Theta^{(2)}[\calX^{2}](x_1,x_2)
%\sum_{q'}\sum_{q} (M_{[1,2],t}(q))_{q'q'} \left( G_{[1,2],t }(x_1) - M_{[1,2],t}(x_1) \right)_{qq}S_{qx_2}
\end{align*}
The first term becomes $\Theta^{(1)}[\calX^{2}](x_1,x_2)$, so the claim follows.
\end{proof}
\end{document}
